% 277

% \S3. Теорема Эйлера.

Пусть $P$ --- произвольный выпуклый многогранник, $e$ --- число его вершин,
$k$ --- чисто рёбер, $f$ -- число граней.
Эйлером\footnote{Леонард Эйлер (1707-1783) великий математик, физик и~астроном.
Швейцарец по происхождению. Был членом Петербургской академии наук. Работал
в~России в~1727-1741 и~в~1766-1783 гг.}
была доказана замечательная теорема.

\begin{Th}
Теорема Эйлера. Для любого выпуклого многогранника

\begin{equation}\label{eq:8_3_1}
e - k + f = 2.
\end{equation}

Соотношение \ref{eq:8_3_1} называется формулой Эйлера.
\end{Th}

Прежде, чем перейти к~доказательству теоремы Эйлера, сделаем несколько замечаний.

В формулировке теоремы Эйлера участвуют лишь числа $e$, $k$ и~$f$, относящиеся
к~элементам выпуклого многогранника --- вершинам, ребрам, граням.
Поэтому можно обращаться к~многогранной поверхности, точнее, к~числу её вершин,
ребер и~граней.

Представим себе, что наша многогранная поверхность деформируется, например,
в~сферу (рис.\ \ref{fig:8_3_20} ).

\begin{figure}\label{fig:8_3_20}
% стр 277 рис 20
\caption{Рис. 20 Деформация поверхности выпуклого многогранника в~сферу.
Её можно представить так --- пусть поверхность многогранника сделана из гибкого
и~растяжимого материала. Затем этот многогранник раздувается --- в~него
накачивается воздух -- и~постепенно деформируется в~сферу. Сеть вершин и~рёбер
многогранника переходит в~сеть на сфере.}
\end{figure}

% 278

При такой деформации сеть вершин и~ребер многогранника перейдёт в~некоторую
криволинейную сеть на сфере. Узлы этой сети естественно назвать вершинами сети,
линии, соединяющие вершины --- рёбрами сети. Части сферы, в~которые перейдут
грани многогранника, условно можно называть гранями сети. Так как при
деформации многогранника числа $e$, $k$ и~$f$ не меняются, то для полученной
на сфере сети также справедлива формула Эйлера:

\begin{equation*}
e - k +f = 2.
\end{equation*}

Вот возникает какой вопрос. На сфере можно представить себе сеть из вершин,
рёбер и~условных граней.
При этом любое ребро соединяет какие-то две вершины сети и~к~любому ребру
прилегают две грани. Кроме того, каждая вершина соединена ребром с~какой-нибудь
другой вершиной. Обозначим $e$, $k$ и~$f$ соответственно число вершин, рёбер
и~граней сети. Будет ли для такой сети на сфере справедлива формула Эйлера
(\ref{eq:8_3_1})? Совершенно ясно, если есть замкнутый выпуклый многогранник,
у которого сеть вершин, ребер и~граней имеет такую же структуру, как и~заданная
сеть на’ сфере, то формуле Эйлера для сети на сфере будет справедлива.

Ответ на вопрос о~существовании выпуклого многогранника с такой же структурой
сети вершин, ребер.и~граней, как и у --- заданной на сфере, лает теорема
Штейница\footnote{Эрнст Штейниц (1871-1928) --- немецкий математик.},
утверждающая, что для каждой сети на сфере (удовлетворяющей
сформулированным выше условиям) можно указать выпуклый замкнутый многогранник,
у~которого сеть вершин, ребер и~граней будет иметь такую же структуру,
что и~заданная на сфере.

Теорему Штейница мы доказывать не будем. Опираясь на эту теорему и~предполагая
доказанной теорему Эйлера, мы можем теперь утверждать, что для любой сети
на сфере справедлива формула Эйлера (\ref{eq:8_3_1}).

При доказательстве теоремы Эйлера мы будем пользоваться наглядно ясным
фактом --- при деформации многогранника число Эйлера $e -k + f$ не меняется.
Эрнст Штейниц (1871-1928) - немецкий математик.

% 279

Доказательство теоремы Эйлера.

Разобьем доказательство на несколько этапов.

1. Пусть $P$ --- выпуклый многогранник, $e$, $k$ и~$f$ --- числа его вершин,
рёбер и~граней. Возьмем какую-либо грань $Q$ многогранника Р (на рисунке
\ref{fig:8_3_21} эта грань $ABCDE$.

\begin{figure}\label{fig:8_3_21}
% стр 279 рис 21
\caption{Рис. 21}
\end{figure}

Затем спроектируем $P$ грань $Q$ из точки $O$ ‚полученной малым смешением
некоторой точки на грани $Q$ во внешнюю область многогранника $P$.
Таким образом мы получим многоугольник (грань $Q$), разбитый некоторой сетью,
состоящей из вершин и~соединяющих их отрезков на простые многоугольники
$T_{1}, T_{2}, \dots T_{f^{\prime}}$\footnote{
Простой многоугольник --- многоугольник, граница которого состоит из одной
замкнутой ломаной.}
(рис.\ \ref{fig:8_3_22}).

\begin{figure}\label{fig:8_3_22}
% стр 279 рис 22
\caption{Рис. 22 Грань $Q$, на которую спроектирован из точки $Q$
(см.\ рис.\ \ref{fig:8_3_21}) многогранник $P$, представляет собой многоугольник,
разбитый на многоугольники $T_{1}, T_{2}, \dots T_{f^{\prime}}$.}
\end{figure}

% 280

2. Подсчитаем число вершин, ребер и~многоугольников только что полученной
сети на грани $Q$.

Число вершин этой сети равно $e$ --- числу вершин многогранника $P$.
Число ребер равно $f$. Число же $f^{\prime}$ многоугольников
$T_{1}, T_{2} \dots T_{f^{\prime}}$‚ равно $f - 1$. где число $f$ --- число
граней многогранника $P$. Поэтому число Эйлера $e - k + f^{\prime}$ для
нашей сети на грани $Q$ будет на единицу меньше числа Эйлера для $P$.
Если мы докажем, что $e - k + f^{\prime}$, то, очевидно, будет
доказана и~теорема Эйлера.

Среди простых многоугольников, на которые разбит многоугольник $Q$, всегда
найдется такой многоугольник (например, многоугольник $T_{1}$, на рисунке
\ref{fig:8_3_22}), что, удалив его из $Q$ мы снова получим один простой
многоугольник $Q$. Если, например, мы удалим из $Q$ многоугольник $T_{1}$‚
изображенный на рисунке \ref{fig:8_3_22}, то оставшийся многоугольник $Q_{1}$,
т.е.\ многоугольник $ABCH^{\prime}KDE$, будет, очевидно, простым. (Попробуйте
обосновать существование многоугольника $T_{1}$ в~общем случае. Вообще говоря,
не каждый из многоугольников $T_{1} \dots T_{f^{\prime}}$‚ выходящих
на границу многоугольника $Q$, обладает таким свойством. На рисунке \ref{fig:8_3_23}
изображён многоугольник $Q$ и~его разбиение на простые многоугольники,
из которых $T_{2}$. не обладает нужным нам свойством).

\begin{figure}\label{fig:8_3_23}
% стр 180 рис 23
\caption{Рис. 23}
\end{figure}

Простой многоугольник $Q$ разбит на простые многоугольники
$T_{1} \dots T_{f^{\prime}}$, Если мы удалим из $Q$ многоугольник $Q_{1}$‚
то останется простой многоугольник $Q_{1}$. Если мы удалим $T_{2}$, то вместо
одного многоугольника $Q$ получим два, а~может быть, и~больше многоугольников.
Но всегда существует такой многоугольник $T_{1}$‚ после удаления которого
получим простой многоугольник $Q_{1}$.

Когда мы удалим $T_{1}$ из $Q$. то при этом удалим из $Q$ часть рёбер и~вершин.
Это ребра и~вершины, не принадлежавшие другим многогранникам из числа
$T_{2} \dots T_{f^{\prime}}$. Ясно, что если мы удалим $m$ рёбер,
то при этом мы удалим $m - 1$ вершину. Число же оставшихся многоугольников
$f_{1}$ равно $f - 1$. Таким образом, число $e_{1}$ вершин многоугольника
$Q_{1}$, будет равно $e - (m - 1)$, число рёбер $k_{1} = k - m$,
а~число $f^{\prime}_{1}$ многоугольников, на которые разбит многоугольник
$Q_{q}$, будет разно $f^{\prime} - 1$.

Подсчитаем число Эйлера $e_{1} - k_{1} + f^{\prime}_{1}$, для $Q_{1}$.
Имеем

\begin{equation*}
  e_{1} - k_{1} + f^{\prime} =
(e - m + 1) + (k - m) + (f^{\prime}) =
e - k + f^{\prime}.
\end{equation*}

Таким образом, при операции удаления из $Q$ многоугольника $f_{1}$
мы получаем многоугольник $Q_{1}$ с~тем же числом Эйлера. Проделав такие
операции $n = f^{\prime} - 1$ раз, мы придём к~одному простому многоугольнику,
для которого число вершин $e_{n}$ равно числу рёбер
$k_{n}$, а~$f^{\prime}_n = 1$.
Поскольку, очевидно, $e_{n} - k_{n} + f^{\prime}_{n} = 1$,
а~$e - k + f^{\prime} = e_{n} - k_{n} + f^{\prime}_{n} = 1$, то равенство
$e - k + f = 1$ справедливо.
Как отмечалось выше, отсюда следует и~справедливость соотношения $e - k + f = 2$.

% 281

\subsubsection{Дополнение. Развёртка выпуклого многогранника.}

Развёрткой называется совокупность простых многоугольников с~указанием правила
склеивания их по сторонам..При этом склеивание двух отрезков означает,
что между их точками устанавливается взаимно однозначное соответствие,
и~соответствующие точки считаются за одну точку.
Предполагается, что выполнены следующие условия.

\begin{enumerate}\label{lst:8_3_1}
\item Склеиваемые отрезки имеют равные длины.
\item Каждая сторона любого многоугольника развертки является стороной одного
многоугольника.\\
Два многоугольника, имеющих общую сторону, называются смежными.
\item Любые два многоугольника $P$ и~$Q$ развёртки конечной цепочкой
многоугольников, в~которой каждый предыдущий многоугольник смежный, с~последующим,
причем первый элемент этой цепочки --- многоугольник $P$.
\item Если многоугольники $P$ и~$Q$ имеют общую вершину, то выбор цепочки,
связывающей $P$ и~$Q$‚ можно оставить так, что все многоугольники этой цепочки
имели общую вершину (рис.\ \ref{fig:8_3_24})

% 282

\begin{figure}\label{fig:8_3_24}
% стр 282 рис 24
\caption{Рис. 24}
\end{figure}

\item Число $e$. вершин, ребер $k$ и~граней (многоугольников) $f$ развёртки
должно удовлетворять формуле Эйлера: $e - k + f = 2$.

Конечно, вершины и~рёбра (стороны) многоугольников, подлежащие склеиванию,
считаются отдельной вершиной и~одним ребром развертки.

\item Сумма плоских углов при каждой из вершин развёртки должна быть меньше
$360^{\circ}$.

\end{enumerate}

Справедлива следующая замечательная теорема, доказанная советским математиком,
академиком А.Д.~Алексанлровым,

\begin{Th}
  Из каждой развертки, удовлетворяющей перечисленным выше условиям \ref{lst:8_3_1},
можно склеить единственный (с~точностью до положения в~пространстве )выпуклый
многогранник.
\end{Th}

Единственность этой теоремы в~более слабой форме быта доказана в~1813~г.
французским математиком Огюстом Коши:

\begin{Th}
Два замкнутых выпуклых многогранника, одинаково составленных из соответственно
равных граней, равны.
\end{Th}

Условие выпуклости в~теореме Коши существенно. Если отказаться от требования
выпуклости, то утверждение теоремы Коши, вообще говоря, не будет справедливым.
Рисунок \ref{fig:8_3_25} поясняет сказанное.

% 283.

\begin{figure}\label{fig:8_3_25}
% стр 283 рис 25
\caption{Рис. 25.}
\end{figure}

В~заключение попытайтесь самостоятельно ответить на следующий вопрос.
На рисунке \ref{fig:8_3_26} изображен многогранник, для которого число Эйлера
$e - k + f$ равно нулю. При этом каждая грань --- простой многоугольник.
Нарисуйте многогранник, для которого число Эйлера $e - k + f$ отлично от 2
и~от 0.

\begin{figure}\label{fig:8_3_26}
% стр 283 рис 26
\caption{Рис. 26}
\end{figure}

Для этого многогранника $e = 12$, $k = 24$, $f = 12$ и~поэтому $e - k + f = 0$.
